\section{Discussion and limitations}\label{sec:discussion}
The six studies answer the four research questions as follows. On RQ1, the spread of the finish-time error grows
almost linearly with the lead time, and a chance constraint built on it reduces to one effective deadline per task that
fast heuristics use unchanged. On RQ2, the planning model matters most, the trigger second and the stopping rule
least. On RQ3, the replanner served more reward than a constant buffer, the nominal model, a conventional planner and
CBBA, and as much as a tuned proportional buffer and moment propagation, but software-in-the-loop flights confirmed only
the fallback. On RQ4, the probability of meeting the deadline decided which given-up tasks to attempt without tuning,
and this beat competitors tuned for the cost of lateness at one task reward, though not at twice that cost.

The order of priorities holds across the designs. A margin that grows with the lead time raised the utility by more
than any choice of stopping rule or trigger, on three independent sets of held-out seeds, without search and with CBBA
in place of ALNS. Periodic replanning helped every method on the nominal model, but only by \val{KDabsAblTrig}\,pp on
the chance-constrained model, which suggests that both correct the same accumulated deviations, one before and one
after they occur. The stopping rule was equivalent to a matched fixed budget in study~2 and added \val{XDabsFullCC}\,pp
in study~4. Its main value lies in avoiding the losses of long budgets on slow planners, probably because the
improvement curve of this ALNS flattens within a few tens of milliseconds.

The shape of the margin matters more than the way it is computed. The constant buffer was clearly worse than every
margin that grows with the lead time, whereas the calibrated effective deadline, a tuned proportional buffer and moment
propagation were within \val{XRangeScaled}\,pp of each other in utility. Because the spread grows almost linearly, a
margin $c\,h$ with a suitable slope captures most of it, and under high disturbances the tuned slope nearly equalled
the one that the calibrated model implies. Moment propagation evaluates the actual legs of a route and reduced the
search rate, but with a cost of lateness its detail paid off in study~6, by \val{WDabsCAMPCAOne}\,pp at one task reward
and \val{WDabsCAMPCATwo}\,pp at two. A planner with a tight budget or a large fleet may therefore prefer the closed
form, and a planner with a generous budget and costly lateness the more detailed model. The comparison with CBBA places
the gain in the planning model rather than in the allocation method. On the proposed model, CBBA came within
\val{WDabsCACBBACAOne}\,pp of the proposed replanner while it planned for \val{WMsCBBACA}\,ms per call, so the
effective deadlines can be handed to an auction-based allocator without changing its bidding rule.

The probability of meeting the deadline is where the calibrated model earns its place. A hard chance constraint
discards every task whose probability falls below $\Phi(z)$, even when a vehicle has nothing better to do, and study~4
shows that this holds however the constraint is evaluated. The nominal fallback recovers these tasks but serves many of
them late. With utility alone, a tuned proportional buffer was as good as the calibrated model. Once a late arrival has
a cost, the question is no longer only how large the margin is but which tasks are worth attempting at all, and that
requires a probability. The closed form turns the break-even rule of Eq.~\eqref{eq:theta} into another effective
deadline, which served more than a buffer and a safety factor tuned for one task reward and stayed ahead of the buffer
whenever the buffer had been tuned for a different cost. Applied to moment propagation, the same rule behaved in the
same way, so it is not tied to the closed form. At twice the cost, a tuned buffer without fallback was better. The rule
prices each attempt on its own and ignores the delay that a fallback task causes to later tasks on its route, and a
route-level version that also charges this delay is the natural next step.

The two tiers play different roles. The kinematic tier resolves differences of a few percentage points over tens of
thousands of missions, whereas a flight with autopilot and vehicle dynamics is costly. Tier B predicted every tier C
comparison on the same missions, so the flights support the transfer of the kinematic results and confirm the
fallback, but the advantage over the nominal model rests on tier B, and tier C studies should be sized from a
conservative effect.

Several limitations bound these conclusions. The error model is linear in the lead time and fitted per disturbance
level, so the level is assumed to be known. Online calibration removes this assumption but served
\val{XDabsOnlCC}\,pp less reward, and its non-inferiority was not shown. The chance constraint is imposed on every task
separately rather than jointly on a route. The fallback, the models of study~4 and the cost-aware fallback were
designed after the held-out results of earlier studies and were therefore tested on new seeds. The cost-aware fallback
assumes that the cost of a late arrival is known and proportional to the reward of the task, and it was tested in
tier B only. CBBA ran centrally with a fully connected network, the most favourable case for its communication, so its
robustness to lost links was not part of the comparison. Tier C uses three UAVs and fifteen tasks per mission, and
software-in-the-loop simulation reproduces the autopilot and the vehicle dynamics but not the sensors, the radio link
or the weather of a field test. The latency factor emulates a slower planning computer instead of measuring one, and
because the stopping rules read a wall clock, the mean utility of a method varied by up to \val{JRepeatDiffMax}\,pp
between two runs of the same missions. Finally, all methods share one family of scenarios, and other mission types may
change the size of the gains.
